\section{Kategorie addytywne i abelowe}\label{cat:sec:abelian-categories}
Punktem wyjścia są grupy abelowe i moduły nad pierścieniem z jedynką (por. \atlasref{AA-16::aa:def:modules-sub-hom}{AA-16: podmoduły i homomorfizmy} oraz \atlasref{AA-16::aa:def:modules-exact}{AA-16: ciągi dokładne}). Chcemy przenieść ich własności na kategorie, w których obiekty nie muszą być zbiorami z wyróżnionymi elementami. Wszystkie kategorie w tym rozdziale są lokalnie małe.

\subsection{Kategorie preaddytywne i biprodukty}
\begin{definicja}[Kategoria preaddytywna]\label{cat:def:preadditive}
Kategoria $\mathcal A$ jest preaddytywna, jeżeli każdy zbiór $\Hom(X,Y)$ jest grupą abelową, a składanie jest addytywne w każdej zmiennej. W szczególności $g0=0$ i $0f=0$: na przykład $g0=g(0+0)=g0+g0$, skąd po odjęciu $g0$ otrzymujemy żądaną równość. Funktor pomiędzy kategoriami preaddytywnymi jest addytywny, jeżeli jego działanie na każdym zbiorze morfizmów jest homomorfizmem grup abelowych.
\end{definicja}
\begin{przyklad}
Kategorie $\Ab$, $\Vect_k$ oraz lewych modułów nad pierścieniem $R$ są preaddytywne: dodawanie morfizmów jest punktowe, a obustronna addytywność składania wynika z addytywności homomorfizmów. Dla modułów por. \atlasref{AA-16::aa:st:modules-hom-abelian}{AA-16: grupa homomorfizmów}. Kategoria grup nieabelowych nie otrzymuje takiej struktury przez punktowe dodawanie: punktowy iloczyn dwóch homomorfizmów nie musi być homomorfizmem.
\end{przyklad}
\begin{definicja}[Biprodukt]\label{cat:def:biproduct}
W kategorii preaddytywnej biproduktem obiektów $X_1,\ldots,X_n$ nazywamy obiekt $B$ z morfizmami $i_j:X_j\to B$, $p_j:B\to X_j$ takimi, że
\[
 p_j i_k=\begin{cases}\id_{X_j}&j=k,\\0&j\ne k,\end{cases}
 \qquad \sum_{j=1}^n i_jp_j=\id_B.
\]
Dla pustej rodziny warunek przyjmuje postać $\id_B=0_{B,B}$.
\end{definicja}
\begin{twierdzenie}[Biprodukt jako iloczyn i koprodukt]\label{cat:thm:biproduct-universal}
Biprodukt jest jednocześnie iloczynem dla projekcji $p_j$ i koproduktem dla włożeń $i_j$. Odwrotnie, każdy skończony iloczyn w kategorii preaddytywnej ma kanoniczną strukturę biproduktu, podobnie każdy skończony koprodukt.
\end{twierdzenie}
\begin{proof}
Dla rodziny $u_j:T\to X_j$ definiujemy $u=\sum_j i_ju_j$. Wówczas $p_ku=\sum_j p_ki_ju_j=u_k$. Jeżeli $v:T\to B$ spełnia te same równania, to $v=\sum_j i_jp_jv=\sum_j i_ju_j=u$. Dla rodziny $v_j:X_j\to T$ definiujemy $v=\sum_j v_jp_j$ i otrzymujemy $vi_k=v_k$. Każde inne $w$ o tych samych składowych spełnia $w=\sum_j wi_jp_j=v$. To dowodzi obu własności uniwersalnych.
Odwrotnie, dla iloczynu $(B,p_j)$ definiujemy $i_k$ jednoznacznie przez warunki $p_ji_k=\delta_{jk}\id$ (dla $j\ne k$ oznacza to morfizm zerowy odpowiedniego typu). Ponieważ $p_k\sum_j i_jp_j=p_k$, jednoznaczność własności uniwersalnej daje $\sum_j i_jp_j=\id_B$. Dla koproduktu dowód otrzymujemy w kategorii przeciwnej, która jest preaddytywna z tymi samymi działaniami na odwróconych zbiorach morfizmów.
Dla $n=0$ warunek $\id_B=0$ pociąga $f=f\id_B=0$ i $g=\id_Bg=0$ dla każdego morfizmu z $B$ i do $B$. Zatem $B$ jest obiektem zerowym. Odwrotnie, obiekt początkowy lub końcowy ma dokładnie jeden endomorfizm, więc jego identyczność jest zerowa.
\end{proof}
\begin{definicja}[Kategoria addytywna]\label{cat:def:additive}
Kategoria addytywna jest kategorią preaddytywną posiadającą wszystkie skończone biprodukty, w tym obiekt zerowy. Biprodukt dwóch obiektów oznaczamy przez $X\oplus Y$.
\end{definicja}
\begin{stwierdzenie}[Rachunek macierzowy]\label{cat:prop:additive-matrices}
Dla biproduktów $X=\bigoplus_{j=1}^nX_j$, $Y=\bigoplus_{i=1}^mY_i$ morfizm $f:X\to Y$ jest jednoznacznie określony przez składowe $f_{ij}=p_i^Yfi_j^X$. Każda rodzina takich składowych wyznacza morfizm $f=\sum_{i,j}i_i^Y f_{ij}p_j^X$. Składanie odpowiada mnożeniu macierzy, w którym mnożenie wpisów oznacza składanie morfizmów.
\end{stwierdzenie}
\begin{proof}
Dla podanego wzoru $p_k^Yfi_l^X=\sum_{i,j}p_k^Yi_i^Y f_{ij}p_j^Xi_l^X=f_{kl}$. Z drugiej strony, wstawiając rozkłady obu identyczności, otrzymujemy
$f=\id_Y f\id_X=\sum_{i,j}i_i^Y(p_i^Yfi_j^X)p_j^X$.
Jeżeli $g:Y\to Z$, to
\[
 (gf)_{kj}=p_k^Zg\left(\sum_i i_i^Yp_i^Y\right)fi_j^X
 =\sum_i g_{ki}f_{ij}.
\]
Dodawanie jest składowe dzięki addytywności składania. Puste sumy są zerami w odpowiednich grupach morfizmów.
\end{proof}

\subsection{Jądra, kojądra, obraz i koobraz}
\begin{definicja}[Jądro i kojądro]\label{cat:def:kernel-cokernel}
Dla $f:X\to Y$ jądrem $\ker f:K\to X$ jest equalizator $f,0$, a kojądrem $\operatorname{coker}f:Y\to Q$ jest koequalizator $f,0$ (por. definicje~\ref{def:equalizer} i \ref{def:coeq}). W szczególności $\ker f$ jest monomorfizmem: jeśli $ku=kv$, to $u,v$ są dwoma rozłożeniami tego samego morfizmu przez jądro, więc $u=v$. Dualnie kojądro jest epimorfizmem. Jądra i kojądra są wyznaczone z dokładnością do jednoznacznego izomorfizmu zgodnego z morfizmem do $X$ lub z $Y$.
\end{definicja}
\begin{stwierdzenie}\label{cat:prop:mono-epi-zero-kernel}
W kategorii addytywnej posiadającej jądra i kojądra morfizm $f:X\to Y$ jest monomorfizmem wtedy i tylko wtedy, gdy jego jądro ma obiekt zerowy jako dziedzinę. Dualnie $f$ jest epimorfizmem wtedy i tylko wtedy, gdy jego kojądro ma obiekt zerowy jako przeciwdziedzinę.
\end{stwierdzenie}
\begin{proof}
Jeżeli $f$ jest monomorfizmem, to $fu=0=f0$ wymusza $u=0$. Morfizm $0\to X$ spełnia więc własność uniwersalną jądra. Odwrotnie, gdy jądro jest zerowe i $fu=fv$, addytywność daje $f(u-v)=0$. Morfizm $u-v$ rozkłada się przez obiekt zerowy, zatem $u-v=0$ i $u=v$. Drugi wniosek wynika z zastosowania pierwszego do kategorii przeciwnej.
\end{proof}
\begin{definicja}[Obraz i koobraz]\label{cat:def:image-coimage}
Zakładamy istnienie wszystkich jąder i kojąder. Dla $f:X\to Y$ przyjmujemy
\[
 k:K\to X=\ker f,\qquad q:Y\to Q=\operatorname{coker}f,
 \qquad p:X\to C=\operatorname{coker}k,\qquad j:J\to Y=\ker q,
\]
gdzie zapis oznacza odpowiednie morfizmy uniwersalne. Obiekty $C,J$ nazywamy odpowiednio koobrazem i obrazem $f$, a morfizmy $p,j$ ich strukturalnymi strzałkami.
\end{definicja}
\begin{stwierdzenie}[Kanoniczne porównanie obrazu i koobrazu]\label{cat:prop:coimage-image}
Istnieje jedyny morfizm $\bar f:C\to J$ taki, że $f=j\bar f p$.
\end{stwierdzenie}
\begin{proof}
Ponieważ $fk=0$, własność kojądra $p$ daje jedyny morfizm $a:C\to Y$ z $ap=f$. Mamy $qap=qf=0$; ponieważ $p$ jest epimorfizmem, $qa=0$. Własność jądra $j$ daje jedyny morfizm $\bar f:C\to J$ z $j\bar f=a$. Jeżeli $jbp=f$, to z epimorficzności $p$ wynika $jb=a$, a z monomorficzności $j$ wynika $b=\bar f$.
\end{proof}
\[
\begin{CD}
X @>{f}>> Y\\
@V{p}VV @AA{j}A\\
C @>{\bar f}>> J
\end{CD}
\]
\begin{definicja}[Kategoria abelowa]\label{cat:def:abelian}
Kategoria abelowa jest kategorią addytywną, w której istnieją wszystkie jądra i kojądra, a kanoniczny morfizm $\bar f$ ze stwierdzenia~\ref{cat:prop:coimage-image} jest izomorfizmem dla każdego $f$.
\end{definicja}
\begin{twierdzenie}[Rozkład morfizmu w kategorii abelowej]\label{cat:thm:abelian-factorization}
W kategorii abelowej każdy morfizm rozkłada się na epimorfizm, po którym następuje monomorfizm. Każdy monomorfizm jest jądrem swojego kojądra, a każdy epimorfizm jest kojądrem swojego jądra. Każdy morfizm będący jednocześnie monomorfizmem i epimorfizmem jest izomorfizmem.
\end{twierdzenie}
\begin{proof}
W rozkładzie $f=j\bar f p$ morfizm $\bar f p$ jest epimorfizmem jako złożenie epimorfizmu i izomorfizmu, a $j$ jest monomorfizmem.
Jeśli $f$ jest monomorfizmem, to $k:0\to X$ jest jego jądrem na mocy stwierdzenia~\ref{cat:prop:mono-epi-zero-kernel}. Identyczność $\id_X$ jest kojądrem $k$, ponieważ każdy morfizm z $X$ zeruje się po złożeniu z $k$ i jednoznacznie rozkłada się przez $\id_X$. Wybierając $p=\id_X$, mamy $f=j\bar f$, gdzie $\bar f$ jest izomorfizmem; stąd $f$ samo spełnia własność uniwersalną jądra $q$: rozłożenie przez $j$ składamy z $\bar f^{-1}$, a jednoznaczność wynika z monomorficzności $f$. Dualnie epimorfizm jest kojądrem swego jądra (kategoria przeciwna kategorii abelowej jest abelowa, gdyż porównanie obrazu i koobrazu staje się odwróconym izomorfizmem).
Jeśli $f$ jest także epimorfizmem, to jego kojądro jest zerowe. Jądrem morfizmu $Y\to0$ jest $\id_Y$, więc $f$ jest izomorfizmem przez jednoznaczność jąder.
\end{proof}
\begin{przyklad}[Moduły]\label{cat:ex:abelian-modules}
W kategorii lewych $R$-modułów biprodukty skończone są sumami prostymi. Jądro $f:M\to N$ jest włączeniem zwykłego jądra, a kojądro jest projekcją $N\to N/f(M)$. Własność uniwersalna jądra wynika z faktu, że $fu=0$ wtedy i tylko wtedy, gdy wartości $u$ należą do jądra; kojądra z własności uniwersalnej modułu ilorazowego. Stąd $C=M/\ker f$ i $J=f(M)$, zaś $\bar f([x])=f(x)$ jest izomorfizmem na mocy \atlasref{AA-16::aa:tw:modules-first-isomorphism}{AA-16: pierwsze twierdzenie o izomorfizmie}. To dowodzi abelowości, w szczególności dla $\Ab$ i $\Vect_k$.
\end{przyklad}
\begin{przyklad}[Kategoria addytywna, lecz nieabelowa]\label{cat:ex:additive-not-abelian}
W pełnej podkategorii skończenie generowanych wolnych grup abelowych istnieją skończone biprodukty. Morfizm $2:\mathbb Z\to\mathbb Z$ (mnożenie przez $2$) jest mono, bo jest iniektywny, i epi: jeśli $u,v:\mathbb Z\to F$ spełniają $u\circ2=v\circ2$, to $2(u(1)-v(1))=0$, skąd $u(1)=v(1)$, ponieważ $F$ jest bez torsji. Nie jest izomorfizmem, gdyż nie jest surjektywny. Stwierdzenie~\ref{cat:thm:abelian-factorization} wyklucza więc abelowość tej kategorii.
\end{przyklad}

\subsection{Ciągi dokładne i rozkładalność}
\begin{definicja}[Dokładność]\label{cat:def:abelian-exact}
W kategorii abelowej ciąg $X\xrightarrow{f}Y\xrightarrow{g}Z$ jest dokładny w $Y$, jeśli $gf=0$ i podobiekty $\operatorname{im}f$ oraz $\ker g$ są równe: istnieje izomorfizm między ich dziedzinami zgodny z monomorfizmami do $Y$. Krótki ciąg dokładny ma postać $0\to A\xrightarrow{i}B\xrightarrow{p}C\to0$, jest dokładny w każdym z trzech niezerowych miejsc (także gdy któryś z obiektów jest zerowy).
\end{definicja}
\begin{stwierdzenie}\label{cat:prop:ses-kernel-cokernel}
Krótki ciąg jest dokładny wtedy i tylko wtedy, gdy $i$ jest jądrem $p$, a $p$ kojądrem $i$.
\end{stwierdzenie}
\begin{proof}
Dokładność na końcach oznacza $\ker i=0$ i $\operatorname{im}p=C$ jako podobiekt $C$. Pierwszy warunek daje monomorficzność $i$. Drugi daje epimorficzność $p$, ponieważ $p$ rozkłada się na epimorfizm i izomorfizm obrazu z $C$. Dla monomorfizmu $i$ jego obraz jest reprezentowany przez samo $i$, więc dokładność w $B$ oznacza, że $i$ jest jądrem $p$. Ponieważ $p$ jest epimorfizmem, jest kojądrem swojego jądra, czyli $i$ (twierdzenie~\ref{cat:thm:abelian-factorization}). Odwrotnie, jądro jest mono, kojądro epi, ich złożenie jest zerowe, a obraz $i$ jest reprezentowany przez $i=\ker p$ jako podobiekt. Epimorfizm $p$ ma obraz całego $C$, gdyż jego kojądro jest zerowe. Wszystkie trzy warunki dokładności są spełnione.
\end{proof}
\begin{twierdzenie}[Rozkładalność krótkiego ciągu dokładnego]\label{cat:thm:abelian-split}
Dla krótkiego ciągu dokładnego $0\to A\xrightarrow{i}B\xrightarrow{p}C\to0$ następujące warunki są równoważne:
\begin{enumerate}
\item istnieje $s:C\to B$ z $ps=\id_C$;
\item istnieje $r:B\to A$ z $ri=\id_A$;
\item istnieje izomorfizm $u:A\oplus C\to B$ taki, że $ui_A=i$ i $pu=p_C$.
\end{enumerate}
\end{twierdzenie}
\begin{proof}
Z (1) wynika $p(\id_B-sp)=0$, więc własność jądra $i$ daje jedyny $r:B\to A$ z $ir=\id_B-sp$. Wówczas $iri=i-spi=i$, czyli $ri=\id_A$, oraz $irs=s-sps=0$, czyli $rs=0$. Obie równości wynikają z monomorficzności $i$. Morfizm $u=[i,s]:A\oplus C\to B$ ma odwrotność $v=\langle r,p\rangle:B\to A\oplus C$: iloczyn $vu$ ma macierz diagonalną z identycznościami, a $uv=ir+sp=\id_B$. Ponadto $ui_A=i$ oraz $pu=p_C$. To daje (3).
Z (2) wynika $(\id_B-ir)i=0$. Własność kojądra $p$ daje jedyny $s:C\to B$ z $sp=\id_B-ir$. Mamy $psp=p$, więc $ps=\id_C$ przez epimorficzność $p$. To daje (1).
Jeśli zachodzi (3), przyjmujemy $s=ui_C$ i $r=p_Au^{-1}$. Wtedy $ps=p_Ci_C=\id_C$ i $ri=p_Au^{-1}i=p_Ai_A=\id_A$.
\end{proof}
\begin{przyklad}
Ciąg $0\to\mathbb Z\xrightarrow{2}\mathbb Z\to\mathbb Z/2\mathbb Z\to0$ jest dokładny, ale nie jest rozkładalny: każdy homomorfizm $\mathbb Z/2\mathbb Z\to\mathbb Z$ jest zerowy, gdyż obraz klasy $1$ musi mieć rząd dzielący $2$. Nie może zatem być prawą odwrotnością projekcji. Por. \atlasref{AA-16::aa:tw:modules-splitting}{AA-16: rozkładalność ciągów modułów}.
\end{przyklad}
Dalsza teoria kompleksów i długich ciągów dokładnych korzysta z tych konstrukcji; w przypadku modułów ich szczegółowe dowody znajdują się w \atlasref{AT-03::def:kompleks-lancuchowy}{AT-03: kompleksy łańcuchowe} oraz \atlasref{AT-03::thm:les-homologia}{AT-03: długi ciąg dokładny homologii}.
